\clearpage
\phantomsection
\addcontentsline{toc}{chapter}{APPENDIX A – CODING}
\begin{center}{\bfseries\fontsize{16}{20}\selectfont APPENDIX A\\[4pt] CODING}\end{center}
\vspace{6pt}
All code is written in Python 3.13 with pandas, scikit-learn, XGBoost, LightGBM, CatBoost, imbalanced-learn, hmmlearn,
SHAP, statsmodels, matplotlib and seaborn. Every notebook is self-contained and runs top to bottom in Google Colab; the
dataset download is included. All randomness is fixed with \texttt{random\_state = 42}, so the numbers in this report
are reproducible. The listings below contain the complete code cells of each notebook.

\renewcommand{\arraystretch}{1.25}
\begin{table}[H]
\centering
\caption{Notebooks and Google Colab links}
\label{tab:colab}
\begin{tabularx}{\linewidth}{@{}l l >{\raggedright\arraybackslash}X@{}}
\toprule
\textbf{Unit} & \textbf{Notebook} & \textbf{Google Colab link} \\
\midrule
1 & \texttt{01\_eda\_play\_store.ipynb} & \colabA \\
1 & \texttt{02\_eda\_video\_game\_sales.ipynb} & \colabB \\
2 & \texttt{03\_regression\_ccpp.ipynb} & \colabC \\
2 & \texttt{04\_classification\_bank\_marketing.ipynb} & \colabD \\
3–4 & \texttt{07\_eda\_har\_accelerometer.ipynb} & \colabE \\
3–4 & \texttt{08\_hmm\_activity\_recognition.ipynb} & \colabF \\
5 & \texttt{09\_diabetes\_ensembles.ipynb} & \colabG \\
\bottomrule
\end{tabularx}
\end{table}

\section*{Chapter 2 – Implementation Methodology and Algorithm Configuration}
\addcontentsline{toc}{section}{Chapter 2 – Implementation Methodology and Algorithm Configuration}
Shared plotting helpers (\texttt{src/viz.py}), followed by the Unit 1 EDA notebooks and the Unit 2 regression and
classification notebooks. Key configurations: 80/20 split, 5-fold (stratified) CV, XGBoost grid search (16
configurations), class weighting / SMOTE, and an F1-optimal threshold chosen on out-of-fold predictions.

\subsection*{A.1 Shared helpers – \texttt{src/viz.py}}
\lstinputlisting{code/viz.py}
\subsection*{A.2 Unit 1 – Play Store EDA (\texttt{01\_eda\_play\_store})}
\lstinputlisting{code/01_eda_play_store.py}
\subsection*{A.3 Unit 1 – Video Game Sales EDA (\texttt{02\_eda\_video\_game\_sales})}
\lstinputlisting{code/02_eda_video_game_sales.py}
\subsection*{A.4 Unit 2 – CCPP Regression (\texttt{03\_regression\_ccpp})}
\lstinputlisting{code/03_regression_ccpp.py}
\subsection*{A.5 Unit 2 – Bank Marketing Classification (\texttt{04\_classification\_bank\_marketing})}
\lstinputlisting{code/04_classification_bank_marketing.py}

\section*{Chapter 3 – Unit 3: Proposed Workflow and Implementation Methodology}
\addcontentsline{toc}{section}{Chapter 3 – Unit 3: Proposed Workflow and Implementation Methodology}
HAR data loading, windowing and feature engineering (notebook 07). The K-Means and GMM clustering with Hungarian
matching and PCA visualisation are in Section 6 of notebook 08 below.
\subsection*{A.6 Units 3–4 – HAR EDA and Features (\texttt{07\_eda\_har\_accelerometer})}
\lstinputlisting{code/07_eda_har_accelerometer.py}

\section*{Chapter 3 – Unit 4: Proposed Workflow and Implementation Methodology}
\addcontentsline{toc}{section}{Chapter 3 – Unit 4: Proposed Workflow and Implementation Methodology}
Unsupervised Gaussian HMM (6 states, full covariance, Baum–Welch, 8 restarts), the hand-coded forward–backward and
Viterbi algorithms with verification, the clustering baselines, and the supervised hybrid HMM (561 features, counted
transition matrix, log-space Viterbi and forward–backward).
\subsection*{A.7 Units 3–4 – HMM, Clustering and Supervised HMM (\texttt{08\_hmm\_activity\_recognition})}
\lstinputlisting{code/08_hmm_activity_recognition.py}

\section*{Chapter 4 – Implementation Methodology and Hyperparameter Configuration}
\addcontentsline{toc}{section}{Chapter 4 – Implementation Methodology and Hyperparameter Configuration}
Diabetes prediction: cleaning, EDA, 2 baselines and 10 ensembles, 5-fold stratified CV, test evaluation, LightGBM
randomised search (40 configurations $\times$ 5 folds, PR-AUC), out-of-fold threshold optimisation, significance tests
and SHAP interpretation. The tuned LightGBM hyperparameters are listed in Table~\ref{tab:tuned}.
\subsection*{A.8 Unit 5 – Ensemble Learning for Diabetes Prediction (\texttt{09\_diabetes\_ensembles})}
\lstinputlisting{code/09_diabetes_ensembles.py}
